\section{The Open Algorithms: What They Need}\label{sec:open}

Twenty of the thirty programs of the folder of algorithms are proved. Of the other ten, eight leave a goal and two
stop because the rewriting exceeds its budget. Their lemmas, written at the beginning of each file, mostly describe
inner loops as families of fresh lists, one per inner trace, with the facts of the inner loop distributed over the
outer iteration constraint; they do not match the closed form of Section~\ref{sec:teq}. For each, this section states
what a proof needs, and, where an attempt was made, what stood in its way.

\subsection{Insertion sort: an inner loop of data-dependent length}\label{sec:open:insertion}
Insertion sort moves the new element down while it is less than its predecessor:
\texttt{while 0 < j \&\& b[j] < b[j - 1] \{ swap; j := j - 1; \}}. Unlike the inner loops of
Section~\ref{sec:sorts}, the number of inner iterations is not determined by the bounds: the constraints over the
inner history state that the comparison held in every inner step, and the exit that it failed or $j$ reached $0$.
The closed form has the inner loop as
\begin{align*}
  \cur(\mathsf{trace}(b,\ &\mathsf{for}\ \bl{\cnt(0, 1, nn)}, \bl{\cnt(i, -1, nn)}\\
  &\mathsf{do}\ x, k', j \mapsto x[j := x[k - k']][k - k' := x[j]]))
\end{align*}
with the position $j - 1$ expressed through the outer iteration $k$, and the constraints are nested inside the outer
iteration constraint, over the history of the outer trace, which the closed form spells out as a term.

\paragraph{What it needs.} A step lemma whose conditions are the inner constraint and exit, and whose conclusion is
that the result extends the sorted prefix by one, used by an induction over the outer trace with the choice among its
dependencies. Both lemmas have conditions that contain inner and outer traces as terms, which runs into the cost
of views recomputed for each instance (Section~\ref{sec:sorts:cycles}); a summary that names the inner loop by
an abstract $\mathsf{insert}(b, i, nn)$ first, as a plain rule, and states the comparisons as facts about the named
result, would keep the traces out of the conditions.

\subsection{The partitions: cycles of three accumulators}\label{sec:open:partition}
In the partition around a pivot, the scanning index $j$ joins the cycle of the array and the boundary $i$, although $j$
is a counter: the guard reads $b[j]$, and every result of a guarded body reads the variables of its guard, including
$j$'s two cases $j + 1$, which are equal. The invariants hold for every run of the update ($0 \le i \le j$, the elements
below $i$ less than the pivot, those from $i$ below $j$ at least the pivot, a permutation), so lemmas in the style of
gnome sort are sound. But the program alone takes 34 s, and 102 s without the aggregation module, since the views of a
cycle of three accumulators are expensive, and with the lemmas the rewriting exceeds its budget.

\paragraph{What it needs.} A read relation that merges equal cases first, or takes the guards only for the results that
differ between the cases, would keep the counter outside the cycle, where the database recognizes it, and leave a
cycle of two. The three-way partition has the same structure with two boundaries.

\subsection{The first duplicate: an early exit by a flag}\label{sec:open:flags}
The first duplicate scans the pairs $i < j$ like the distinctness check, but stops both loops when the flag is set. The
inner flag is the one of Section~\ref{sec:sorts:distinct}, $r \vee \neg\mathsf{nodup}(a, i, i + 1 + \intf(nn))$. The outer
loop runs while the flag is false, so its constraint states that the flag was false at the head of every outer iteration,
over the history of the outer trace; the exit of the outer loop is a disjunction, $n \le i \vee r$, and the exit of the
inner loop depends on the flag as well.

\paragraph{What it needs.} Two rules on the atom $\mathsf{distincts}(\ldots)$, one per case of the outer exit: if the flag is
set, the last iteration found a duplicate within $[0, n)$; if it is not, every iteration ran with its full length and
found none. Both take the constraint on the flag at the heads as a condition, which contains the outer trace as a term;
with the postcondition as two implications, each goal has its literal of the exit.

\subsection{The reversal: a fact lost in a normal form}\label{sec:open:reverse}
The swaps of $k$ and $n - 1 - k$ are recognized as a reversal by a lemma of the file, whose conditions are linear:
$2 n' \le c + 1$ and $c \le 2 n'$ for the start $c = n - 1$ of the counter down. They follow from the facts if
$\intf(n') = i$ for the final value $i$ of the counter. The database rewrites $\intf(n) = x$ into $n = \natf(x)$ under the
condition $0 \le x$, and the result no longer states $0 \le x$, which is lost if it followed from the equation itself.
The equivalence $\intf(n) = x \Leftrightarrow 0 \le x \wedge n = \natf(x)$ holds without a condition and keeps the fact; with it,
the conditions of the reversal follow.

\subsection{Witnesses, duplicates and the step budget}\label{sec:open:other}
\begin{itemize}[nosep]
\item \emph{An existential postcondition} (\texttt{max\_element}): the maximum is
  \[ m = \cur(\mathsf{maxs}(a[0], \mathsf{selects}(a, \bl{\cnt(1, 1, n)}))), \]
  and $\exists k.\ a[k] = m$ needs a witness. A lemma whose left-hand side is the existential does not apply: the
  matcher has no case for quantifiers. The same program with the postcondition $\mathsf{occurs}(m, \ldots)$ in the
  vocabulary of the database closes; quantified left-hand sides, or witnesses taken from the terms of the goal, would
  close the other.
\item \emph{The transpose} of the folder keeps two folds and lemmas for fresh lists; the same program with one lemma
  on the nest (Section~\ref{sec:transpose}) closes.
\item \emph{The step budget}: the cocktail sort, with two inner loops per outer iteration, and the counting sort, with
  four loops in sequence and a scatter through a table of positions, stop because the rewriting exceeds its budget,
  with a rule on counters or sums as the last step. The last step of a rewriting that does not terminate is rarely its
  cause: finding the cycle needs the rules of the file in pairs on the program, or a report of the rule applications
  that repeat.
\end{itemize}
